\chapter{Kako se čita i kako se crta graf}
\label{ch:graf}

U ovom kolegiju slika nije ukras nego dokaz. Kad na obrani kažeš ,,rješenje je u
ovoj točki``, ispitivač će te gotovo sigurno zamoliti da to \emph{pokažeš} --- da
nacrtaš dopustivi skup, nivo-krivulje i gradijente, i da s crteža pročitaš zašto
je baš ta točka optimum.

Ovo poglavlje objašnjava svaki element takve slike: što je, odakle dolazi, kako
ga izračunati i kako ga nacrtati rukom, bez računala.

\section{Koordinatni sustav i što je na kojoj osi}
\label{sec:koordinatni-sustav}

U Zadacima 1 i 2 crta se \textbf{prostor projektnih varijabli}: na vodoravnoj
osi je $x_1$, na okomitoj $x_2$. Svaka točka ravnine je jedan mogući izbor
projektnih varijabli. To nije graf funkcije --- vrijednost funkcije cilja nije
ni na jednoj osi, nego se prikazuje nivo-krivuljama.

U Zadatku 3 crta se \textbf{fizikalni presjek terena}: na vodoravnoj osi je
koordinata $x \in [0, L]$ uzduž trase, na okomitoj visina. Dvije krivulje na
istoj slici su teren $h(x)$ i niveleta ceste $y(x)$.

U Zadatku 4 crtaju se \emph{dvije različite} vrste slika:
\begin{itemize}[nosep,leftmargin=*]
  \item \textbf{vremenski odziv} --- na vodoravnoj osi vrijeme $t$, na okomitoj
        komponente stanja $x_1(t)$ i $x_2(t)$, dvije krivulje;
  \item \textbf{fazni portret} --- na vodoravnoj osi $x_1$, na okomitoj $x_2$,
        a vrijeme je \emph{sakriveno}: vidi se samo putanja koju točka opisuje.
\end{itemize}

\begin{zamka}[Zamka: zamjena vremenskog odziva i faznog portreta]
Zadatak 4b traži oboje, i to su različite slike istog rješenja. Na faznom
portretu nema vremenske osi, pa se s njega \emph{ne} može očitati koliko je
dugo trajala tranzicija --- samo kojim putem je stanje išlo. Obrnuto, iz
vremenskog odziva ne vidi se spiralno namotavanje. Zato se traže obje.
\end{zamka}

\section{Nivo-krivulje}
\label{sec:nivo-krivulje}

\begin{pojam}{Nivo-krivulja (nivo-skup)}
Za funkciju $f: \R^2 \to \R$ i broj $c \in \R$, \textbf{nivo-krivulja razine $c$} je
skup
\begin{equation}
S_c \coloneqq \{\, x \in \R^2 \mid f(x) = c \,\}.
\end{equation}
To je skup svih točaka u kojima funkcija ima \emph{istu} vrijednost. Čita se:
,,es ce``, ili ,,nivo-skup razine ce``.
\end{pojam}

Analogija koja pomaže: na zemljopisnoj karti to su izohipse. Sve točke na istoj
izohipsi su na istoj nadmorskoj visini. Gušće izohipse znače strmiji teren.
Točno to vrijedi i ovdje: gdje su nivo-krivulje gušće, funkcija cilja brže
raste.

\subsection{Zašto su nivo-krivulje kvadratne funkcije elipse}

Za $f(x) = \tfrac{1}{2}x\trans Qx + c\trans x$ s $Q \pd 0$ vrijedi
\begin{equation}
f(x) = \tfrac{1}{2}(x - x^{0})\trans Q (x - x^{0}) + f(x^{0}),
\qquad x^{0} \coloneqq -Q^{-1}c ,
\label{eq:dopuna-kvadrata}
\end{equation}
što je matrična verzija ,,nadopune do potpunog kvadrata``. Uvjet $f(x) = c$ tada
postaje
\begin{equation}
(x - x^{0})\trans Q (x - x^{0}) = 2\bigl(c - f(x^{0})\bigr) = \text{konst.},
\end{equation}
a to je jednadžba \textbf{elipse sa središtem u $x^{0}$}. Za $c < f(x^{0})$ skup
je prazan --- ispod minimuma funkcija ne poprima nijednu vrijednost.

\begin{korak}[Elipse za funkciju cilja iz Zadatka 1]
\begin{enumerate}[leftmargin=*]
  \item \textbf{Središte}: $x^{0} = -Q^{-1}c$. Uz $Q^{-1} = \tfrac{100}{791}\begin{bsmallmatrix}4 & 0{,}3 \\ 0{,}3 & 2\end{bsmallmatrix}$
        i $c = \begin{bsmallmatrix}2 \\ -3\end{bsmallmatrix}$:
        \[
          x^{0} = -\frac{100}{791}\begin{bmatrix}4 & 0{,}3 \\ 0{,}3 & 2\end{bmatrix}
                   \begin{bmatrix}2 \\ -3\end{bmatrix}
                = -\frac{100}{791}\begin{bmatrix}7{,}1 \\ -5{,}4\end{bmatrix}
                = \begin{bmatrix}-710/791 \\ 540/791\end{bmatrix}
                \approx \begin{bmatrix}-0{,}898 \\ 0{,}683\end{bmatrix}.
        \]
        Središte elipsa leži \emph{izvan} dopustivog skupa --- zato optimum
        završi na rubu.
  \item \textbf{Smjerovi osi}: to su svojstveni vektori matrice $Q$. Za
        $\lambda_1 = 1{,}956$ iz $(2 - \lambda_1)v_1 - 0{,}3\,v_2 = 0$ slijedi
        $v_2 = 0{,}1468\,v_1$, dakle smjer pod kutom
        $\arctan(0{,}1468) \approx 8{,}4^\circ$ prema osi $x_1$. Drugi svojstveni
        vektor je na njega okomit, pod $\approx 98{,}4^\circ$.
  \item \textbf{Koja je os duža}: manja svojstvena vrijednost daje \emph{dužu}
        poluos, jer je duž tog smjera funkcija ,,plića``. Omjer poluosi je
        $\sqrt{\lambda_2/\lambda_1} = \sqrt{4{,}044/1{,}956} \approx 1{,}44$.
  \item \textbf{Zaključak za crtanje}: elipse su blago spljoštene (omjer $\approx 3{:}2$),
        blago nagnute (oko $8^\circ$), koncentrične oko točke $(-0{,}9,\; 0{,}68)$.
\end{enumerate}
\end{korak}

\begin{ploca}[Nivo-elipse na ploči, bez računala]
\begin{enumerate}[leftmargin=*]
  \item Nacrtaj osi i označi središte $x^{0} \approx (-0{,}9,\; 0{,}7)$ malim križićem.
  \item Kroz središte povuci tanku iskošenu crtu pod otprilike $8^\circ$ --- to je
        duža os. Okomito na nju drugu tanku crtu --- kraća os.
  \item Nacrtaj dvije-tri koncentrične elipse oko središta, svaku sve veću, s
        omjerom poluosi otprilike $3{:}2$ uzduž tih crta.
  \item Uz svaku napiši rastuću vrijednost, npr.\ $f = 5$, $f = 15$, $f = 30$.
        Ne moraju biti točne; poanta je da rastu prema van.
  \item Reci naglas: ,,elipse rastu prema van, pa minimum tražimo što bliže
        središtu --- a ograničenja nas guraju od njega``.
\end{enumerate}
\end{ploca}

\section{Gradijent na slici}
\label{sec:gradijent-na-slici}

\begin{izvor}{Predavanje 2 --- Definicije, klasifikacija, konveksnost, slajd 37/71 (PDF str.~38)}
\slajd[0.86]{p02-gradijent-okomit}

\textbf{Transkripcija --- dva svojstva gradijenta:}
\begin{enumerate}[nosep,leftmargin=*]
  \item Gradijent funkcije $f$ u točki $x$ daje smjer \emph{najbržeg rasta} funkcije.
  \item Gradijent $\grad f(x)$ je ortogonalan (okomit) na tangentu nivo-krivulje
        $S_c$ u bilo kojoj točki $x$ za koju je $f(x) = c$.
\end{enumerate}

\textbf{Zašto je drugo svojstvo istinito} (profesor to dokazuje na sljedećem
slajdu): parametriziraj nivo-krivulju kao $g(t)$ i definiraj $h(t) = f(g(t))$.
Po definiciji nivo-krivulje $h(t) = c$ za sve $t$, pa je $\mathrm{d}h/\mathrm{d}t = 0$.
Lančano pravilo daje $0 = \grad f(x)\trans d$, gdje je $d$ tangentni vektor
krivulje. Skalarni produkt jednak nuli znači okomitost.
\end{izvor}

Praktična posljedica za crtanje: gradijent nikad ne crtaš ,,po osjećaju``. Nacrtaš
nivo-krivulju kroz točku, pa strelicu \emph{okomito} na nju, usmjerenu prema
većim vrijednostima --- prema van, ako su elipse koncentrične oko minimuma.

\begin{formula}{Gradijenti u Zadatku 1, brojevima}
\begin{align}
\grad f(x) &= Qx + c
  && \text{ovisi o točki, mijenja se po ravnini,} \\
\grad g_1(x) &= \begin{bmatrix} -1 \\ -1 \end{bmatrix}
  && \text{konstantan, jer je } g_1 \text{ afina,} \\
\grad g_2(x) &= \begin{bmatrix} 1 \\ -1 \end{bmatrix}
  && \text{konstantan, jer je } g_2 \text{ afina.}
\end{align}
\end{formula}

\begin{zamka}[Zamka: duljina strelice gradijenta]
Na slici gradijente crtamo skaliranima, da stanu i da se vide. To je dopušteno
--- \emph{smjer} je ono što nosi informaciju --- ali to na obrani moraš reći.
Ako prešutiš, ispitivač te može pitati zašto $\grad f$ izgleda jednako dug kao
$\grad g_1$ iako im norme nisu jednake. Odgovor: ,,strelice su normirane radi
preglednosti, uspoređujem samo smjerove i međusobne kutove``.
\end{zamka}

\section{Dopustivi skup}
\label{sec:dopustivi-skup}

Dopustivi skup je presjek poluprostora --- postupak crtanja svakog poluprostora
razrađen je u \S\ref{sec:poluprostor}. Nakon što ucrtaš oba rubna pravca i
odrediš dopustivu stranu svakoga, dopustivi skup je područje koje zadovoljava
\emph{oba} uvjeta istovremeno; obično se šrafira ili osjenča.

\begin{korak}[Redoslijed crtanja cijele slike Zadatka 1]
\begin{enumerate}[leftmargin=*]
  \item Osi $x_1$, $x_2$ s brojevima; raspon otprilike $0$ do $10$ po obje osi.
  \item Pravac $x_1 + x_2 = 10$ kroz $(0,10)$ i $(10,0)$.
  \item Pravac $x_1 - x_2 = 3$ kroz $(3,0)$ i $(0,-3)$.
  \item Testiraj ishodište u obje nejednakosti i osjenčaj presjek dopustivih
        strana --- ispada neograničeno klinasto područje gore desno.
  \item Nivo-elipse oko $x^{0} \approx (-0{,}9,\; 0{,}7)$, prema
        \S\ref{sec:nivo-krivulje}.
  \item Točku optimuma $\xopt$ označi punim kružićem i napiši joj koordinate.
  \item Iz $\xopt$ nacrtaj tri strelice: $\grad f(\xopt)$, $\grad g_1$, $\grad g_2$.
  \item Provjeri očima: leži li $-\grad f(\xopt)$ unutar stošca razapetog
        gradijentima \emph{aktivnih} ograničenja? Ako da, KKT uvjeti su
        geometrijski potvrđeni.
\end{enumerate}
\end{korak}

\section{Kako sliku pročitati naglas}
\label{sec:citanje-slike}

Ispitivač rijetko pita ,,što je na slici``. Pita ,,što slika \emph{pokazuje}``.
Razlika je u tome što drugo pitanje traži zaključak, ne opis.

\begin{izgovor}[Predložak za svaku sliku, četiri rečenice]
\begin{enumerate}[leftmargin=*]
  \item \textbf{Osi:} ,,Na apscisi je $x_1$, na ordinati $x_2$; ovo je prostor
        projektnih varijabli.``
  \item \textbf{Elementi:} ,,Osjenčano je dopustivo područje, presjek dvaju
        poluprostora; koncentrične elipse su nivo-krivulje funkcije cilja, a
        rastu prema van.``
  \item \textbf{Ključna točka:} ,,Optimum je u ovoj točki na rubu, jer je
        središte elipsa izvan dopustivog područja --- neograničeni minimum nije
        dopustiv.``
  \item \textbf{Zaključak:} ,,U optimumu je nivo-elipsa tangentna na aktivno
        ograničenje, pa su gradijent cilja i gradijent tog ograničenja
        kolinearni --- to je upravo uvjet stacionarnosti iz KKT sustava.``
\end{enumerate}
\end{izgovor}

\begin{pitanje}{Što fizikalno znači to da je nivo-krivulja tangentna na ograničenje u optimumu?}
\textbf{Odgovor.} Da se duž ruba dopustivog skupa vrijednost funkcije cilja u
prvom redu više ne mijenja. Kad bi nivo-krivulja rub \emph{sjekla}, postojao bi
smjer uzduž ruba koji ostaje dopustiv, a smanjuje cilj --- pa točka ne bi bila
optimum. Tangentnost je geometrijski oblik uvjeta da projekcija gradijenta cilja
na dopustive smjerove iščezava, dakle uvjeta stacionarnosti.
\end{pitanje}

\section{Vremenski odziv i fazni portret}
\label{sec:odziv-fazni}

Za Zadatak 4 vrijede druga pravila čitanja.

\begin{itemize}[leftmargin=*]
  \item \textbf{Vremenski odziv}: gledaš \emph{omotnicu}. Ako vrhovi oscilacija
        rastu, sustav divergira; ako padaju prema nuli, konvergira. Kod sklopnog
        sustava korisno je okomitim isprekidanim crtama označiti cijele
        vrijednosti $t$ --- trenutke u kojima se mijenja $A(t)$ --- pa se vidi
        što se događa unutar svakog intervala.
  \item \textbf{Fazni portret}: gledaš \emph{namotava li se} putanja prema
        ishodištu ili se odmotava od njega. Spirala prema unutra znači
        asimptotsku stabilnost, spirala prema van nestabilnost. Početna točka
        $x(0) = (1,\,0)$ uvijek se posebno označi.
\end{itemize}

\begin{izgovor}
,,Na faznom portretu otvorenog kruga putanja se odmotava od ishodišta, dakle
sustav je nestabilan --- iako su, kako sam pokazao pod a), oba pojedinačna moda
Hurwitzova. To je bit ovog zadatka: stabilnost svakog moda posebno ne povlači
stabilnost sklopnog sustava.``
\end{izgovor}
